\cotefile{93}
\scanpage{3}
\margin{\handwritten{G? = correction de Grothendieck}}
\margin{\handwritten{Pour une version un brin plus \unsure{complete|complète} : Lecture Notes 476 (Anvers $\overline{\underline{\mathrm{IV}}}$) p 53--73 see also : Tate, ibidem p 33--52 et \unsure{spécialement} p 33--39}}

\begin{center}
COURBES ELLIPTIQUES : FORMULAIRE

d'après J. Tate\\
mis au goût du jour par P. Deligne
\end{center}

1.\qquad \uline{Formules générales}.

Dans ce formulaire, on appellera courbe de genre 1 sur un schéma de base $S$ un morphisme propre et plat de présentation finie
\[
p : E \handwritten{\longrightarrow} S \qquad ,
\]
muni d'une section $e$ contenue dans l'ouvert de lissité, dont les fibres géométriques sont des courbes réduites irréductibles de genre \margin{\handwritten{G?}}\handwritten{\uline{ar}}ithmétique 1 . Les fibres géométriques de $p$ sont donc soit :

a) une courbe elliptique, i.e. une courbe propre, lisse et connexe de genre 1 ,

b) une droite projective dont on a identifié deux points distincts (cubique dans $P_2$ à point double ordinaire),

c) une droite projective dont on a identifié deux points infiniment voisins (cubique dans $P_2$ à singularité cuspidale).

On dira dans le cas a) (resp. b), c)) que la fibre est de type elliptique (resp. multiplicatif, resp. additif).

Si $E$ est une courbe de genre 1 sur $S$, on définit un faisceau inversible $\omega$ sur $S$ par la formule
\[
\omega=e^*\Omega^1_{E/S}\qquad .
\]
La section $e$ est un diviseur de Cortier relatif de $E$ sur $S$ ;

% archive sheet number: 2
\scanpage{4}
\uline{Formules utiles}

\[
1728\Delta=c_4^3-c_6^2
\]
\[
j=\frac{c_4^3}{\Delta}=1728+c_6^2/\Delta
\]
\[
c_4\in H^0(\unsure{\bar{\Pi}|\bar{M}},\omega^4)\qquad
c_6=H^0(\bar{\Pi},\omega^6)\qquad \Delta=H^0(\bar{\Pi},\omega^{12}
\]
\margin{car $\ne 2,3$ .}
\[
\begin{aligned}
c_4=0&\Longleftrightarrow 6\ \text{auto}\Longleftrightarrow j=0\\
c_6=0&\Longleftrightarrow 4\ \text{auto}\Longleftrightarrow j=1728
\end{aligned}
\]
\[
h\in H^0(\bar{\Pi},\omega^{\otimes\deleted{\unsure{2|(p-1)/2}}\inserted{p-1}})
\]
\[
\deg\omega=1/24
\]
\begin{center}\rule{1.4cm}{0.3pt}\end{center}
\[
\begin{aligned}
E=G_m/q^{\mathbb Z},\ \frac{du}{u}\quad\to\quad
c_4&=\left(\frac{du}{u}\right)^{\otimes4}
\left(1+240\sum\sigma_3(n)q^{\unsure{n}}\right)\quad(=E_4)\quad(Q)\\
c_6&=\left(\frac{du}{u}\right)^{\otimes6}
\left(1\deleted{\unsure{+}}\unsure{-}504\sum\sigma_5(n)q^n\right)\quad(=E_6)\quad(R)
\end{aligned}
\]
\begin{center}\rule{1.4cm}{0.3pt}\end{center}

Question : $h=\left(\dfrac{du}{u}\right)^{\otimes p-1}$ ? à l'$\infty$ -

\begin{center}\rule{1.6cm}{0.3pt}\end{center}
\[
y^2=x(x-1)(x-\lambda)\quad\longmapsto\quad
j=2^8\frac{(1-\lambda(1-\lambda))^3}{(\lambda(1-\lambda))^2}
\]
$[$\unsure{p\textsuperscript{ts}|p\textsuperscript{t}} ordre 2 marqués$]$
\diagram{\(\boxed{\begin{array}{l}
\lambda\\1-\lambda\\\lambda^{-1}\\1-\lambda^{-1}\\(1-\lambda)^{-1}\\
\unsure{-|1-}(1-\lambda^{-1})^{-1}
\end{array}}\)}
\begin{center}\rule{1.8cm}{0.3pt}\end{center}

\uline{Hasse} : $\displaystyle(-1)^r\sum\binom{r}{i}^2\lambda^i$

ou $r=\dfrac{p-1}{2}$

\diagram{\(\begin{array}{rcccccl}
p\equiv&1&5&7&11&\cdot&(12)\\
\lambda=&&0&1728&\substack{0\\1728}&&\text{et }\dfrac{p-\unsure{\cdot|0}}{\ldots12}\text{ autres }-
\end{array}\)}
\begin{center}\rule{1.6cm}{0.3pt}\end{center}
% archive sheet number: 3
\scanpage{5}
\mspage{2}
de plus, on vérifie fibre par fibre que
\[
R^1p_*\handwritten{\mathscr O}(ne)=0\quad\text{pour}\quad n>0\qquad .
\]

Il en résulte, par Riemann-Roch, que pour $n>0$, le faisceau $p_*\handwritten{\mathscr O}(ne)$ est localement libre de rang $n$, et que sa formation commute à tout changement de base. La suite exacte de cohomologie fournit de plus une suite exacte
\[
0\handwritten{\longrightarrow}p_*\handwritten{\mathscr O}(ne)\handwritten{\longrightarrow}
p_*\handwritten{\mathscr O}((n+1)e)\handwritten{\longrightarrow}\omega^{\otimes-(n+1)}
\handwritten{\longrightarrow}0\quad(n>0)\quad .
\]
Par ailleurs, $\handwritten{\mathscr O}(e)$ étant localement libre de rang 1 , on vérifie fibre par fibre, compte tenu \unsure{dé|de} (SGA 1 X 1.2), que
\[
\handwritten{\mathscr O}_S\handwritten{\xrightarrow{\sim}}p_*\handwritten{\mathscr O}_E
\handwritten{\xrightarrow{\sim}}p_*\handwritten{\mathscr O}(e)\qquad .
\]
En d'autres termes, les $p_*\handwritten{\mathscr O}(me)$ pour $1\leq m\leq n$ définissent une filtration de $p_*\handwritten{\mathscr O}(ne)$, de gradué associé
\[
\mathrm{Gr}\ p_*\handwritten{\mathscr O}(ne)=\sum_{\substack{i=0\\i\ne1}}^n\omega^{\otimes-i}\qquad .\tag*{(1.1)}
\]

Le morphisme $p$ est d'intersection complète relative ; il en résulte que $Rp^!\handwritten{\mathscr O}_S\in D^+(E)$ \handwritten{\tikz[baseline=-.2ex,x=1ex,y=1ex]{\draw (0,0)--(.28,1.22) .. controls (.13,1.48) and (.51,1.55) .. (.47,1.35) .. controls (.44,1.22) and (.30,1.22) .. (.28,1.22);}}\margin{\handwritten{\uline{réf}}} a un seul groupe de cohomologie non nulle, $R^1p^!\handwritten{\mathscr O}_S$, et que ce dernier est un faisceau inversible sur $E$ ; on l'appelle le faisceau des différentielles régulières sur $E$ et on le désigne par $\Omega^{\mathrm{reg}}_{E/S}$. Là où $p$ est lisse, on a $\Omega^{\mathrm{reg}}_{E/S}=\Omega^1_{E/S}$. La dualité montre que $R^1p_*\Omega^{\mathrm{reg}}_{E/S}$ est isomorphe à $\handwritten{\mathscr O}$, $p_*\Omega^{\mathrm{reg}}_{E/S}$ est donc\handwritten{\tikz[baseline=-.2ex,x=1ex,y=1ex]{\draw (0,0)--(.28,1.22) .. controls (.13,1.48) and (.51,1.55) .. (.47,1.35) .. controls (.44,1.22) and (.30,1.22) .. (.28,1.22);}}\margin{\handwritten{\uline{réf}}} localement libre, de formation compatible à tout changement de base et, raisonnant fibre par fibre, on en déduit que

% archive sheet number: 4
\scanpage{6}
\mspage{3}
\[
p_*\Omega^{\mathrm{reg}}_{E/S}\handwritten{\xrightarrow{\sim}}o^*\Omega^1_{E/S}=\omega ,
\]
et que pour la flèche adjointe à la flèche inverse
\[
p^*\omega\handwritten{\xrightarrow{\sim}}\Omega^{\mathrm{reg}}_{E/S}\qquad .
\]
Les sections de $\Omega^{\mathrm{reg}}_{E/S}$ ainsi associées aux sections de $\omega$ s'appellent les différentielles invariantes.

Soit $\pi$ une section inversible de $\omega$ ; il existe, localement pour la topologie de Zariski, une base $(y,x,1)$ de $p_*\handwritten{\mathscr O}(3e)$ telle que l'image de $y$ dans $\omega^{\otimes-3}$ soit $-\pi^{\otimes-3}$, que $x$ appartienne à \margin{\handwritten{G?}}$p_*\handwritten{\mathscr O}(2e)$, et que son image dans $\omega^{\otimes-2}$ soit $\pi^{\deleted{\unsure{8|0|\otimes}}\otimes-2}$. Si $\pi'$ est une autre section inversible de $\omega$, les bases correspondantes s'écrivent toutes, sous forme implicite (pour $r,s,t$ convenables)
\[
\handwritten{\left\{}\begin{aligned}
x&=u^2x'+r\\
y&=u^3y'+su^2x'+t\\
\pi'&=u\pi\qquad .
\end{aligned}\right.\tag*{(1.2)}
\]

Le faisceau\unsure{\textbackslash|'} $\handwritten{\mathscr O}(3e)$ est très ample et $(x,y,1)$ sont les coordonnées projectives d'un plongement de $E$ dans $P_2/S$.

La courbe image vérifiera une équation
\[
y^2+a_1xy+a_3y=x^3+a_2x^2+a_4x+a_6\qquad .\tag*{(1.3)}
\]

Pour vérifier l'existence et l'unicité de cette équation, on

% archive sheet number: 5
\scanpage{7}
\mspage{4}
utilise (1.1), en regardant les deux membres de (1.3) comme des sections de $p_*\handwritten{\mathscr O}(6e)$ : en effet, $y^2-x^3$ est dans $p_*\handwritten{\mathscr O}(5e)$, $xy$ dans $p_*\handwritten{\mathscr O}(5e)$, $x^2$ dans $p_*\handwritten{\mathscr O}(4e)$, $y$ dans $p_*\handwritten{\mathscr O}(3e)$, $x$ dans $p_*\handwritten{\mathscr O}(2e)$ et $1$ dans $p_*\handwritten{\mathscr O}$.

Si on pose
\[
F(x,y)=y^2+a_1xy+a_3y-x^3-a_2x^2-a_4x-a_6
\]
la forme différentielle invariante définie par $\pi$ s'écrit
\[
\begin{aligned}
\pi&=dx/F_y=-dy/F_x\qquad\text{soit}\\
\pi&=\frac{dx}{2y+a_1x+a_3}=\frac{dy}{3x^2+2a_2x+a_4-a_1y}\qquad .
\end{aligned}
\]

Réciproquement toute équation (1.3) est l'équation non homogène d'une courbe de genre 1 plongée comme une cubique dans $P_2$ et présentant un point d'inflexion en $(0,1,0)$ avec la droite de l'infini pour tangente d'inflexion.

On pose alors
\[
b_2=a_1^2+4a_2\qquad b_4=a_1a_3+2a_4\qquad b_6=a_3^2+4a_6\tag*{(1.4)}
\]
\[
-b_8=a_1a_3a_4+a_4^2-a_1^2a_6-a_2a_3^2-4a_2a_6
\]
\[
d_6=b_2b_4-18b_6\qquad d_8=b_2^2b_4-2^4\,3\,b_4^2+2\,3^2\,b_2b_6
\]
On a
\[
\begin{aligned}
&-4b_8=b_4^2-b_2b_6\\
&12d_6=c_6+b_2c_4\qquad -12d_8=c_4^2+b_2c_6
\end{aligned}
\]
% archive sheet number: 6
\scanpage{8}
\mspage{5}
\[
c_4=b_2^2-24b_4\qquad -c_6=b_2^3-36b_2b_4+216b_6\tag*{(1.5)}
\]
\[
\begin{aligned}
\Delta&=-b_2^2b_8-8b_4^3-27b_6^2+9b_2b_4b_6\\
&=b_4^3-27b_6^2+b_8(36b_4-b_2^{\handwritten{2}})\qquad .
\end{aligned}
\]

Soient maintenant $(x',y',\pi')$ reliés à $(x,y,\pi)$ par (1.2). On aura :
\[
\begin{aligned}
ua'_1&=a_1+2s\\
u^2a'_2&=a_2-sa_1+3r-s^2\\
u^3a'_3&=a_3+ra_1+2t=F_t(r,t)\\
u^4a'_4&=a_4-sa_3+2a_2r-(t+rs)\unsure{{}'|\cdot}a_1+3r^2-2st=-F(r,t)-sF_t(r,t)\\
u^6a'_6&=a_6+ra_4+r^2a_2+r^3-ta_3-t^2-rta_1=-F(r,t)
\end{aligned}\tag*{(1.6)}
\]
\[
\begin{aligned}
u^2b'_2&=b_2+12r\\
u^4b'_4&=b_4+rb_2+6r^2\qquad(\text{équ. de discr. }c_4)\\
u^6b'_6&=b_6+2rb_4+r^2b_2+4r^3\qquad(\text{équ. de discr. }16\,\Delta)\\
u^8b'_8&=b_8+3rb_6+3r^2b_4+r^3b_2+3r^4\\
u^6d'_6&=d_6+c_4r\qquad u^8d'_8=d_8-c_6r
\end{aligned}\tag*{(1.7)}
\]
% archive sheet number: 7
\scanpage{9}
\mspage{6}
\[
u^4c'_4=c_4\qquad u^6c'_6=c_6\qquad u^{12}\Delta'=\Delta\qquad .\tag*{(1.8)}
\]

Les formules (1.8) expriment que la section $c_4\pi^{\otimes4}$ (resp. $c_6\pi^{\otimes6}$, resp. $\Delta\pi^{\otimes12}$) du fibré $\omega^{\unsure{\tikz[baseline=.06ex,x=1ex,y=1ex]{\draw[line width=.055ex] (.4,.4) circle (.33); \draw[line width=.055ex] (.17,.63)--(.63,.17);}|\otimes}4}$ (resp. $\omega^{\otimes6}$, $\omega^{\bigcirc12}$) ne dépend pas du choix arbitraire de $\pi$, $x$ et $y$.

2.\qquad \uline{Cas où 2 et 3 sont inversibles}.

Lorsque 2 et 3 sont inversibles, on peut d'une et d'une seule façon choisir $x$ et $y$ ($\pi$ étant donné) de sorte que dans l'équation (1.3.) on ait $a_1=a_2=a_3=0$. On a alors
\[
y^2=x^3+a_4x+a_6\qquad,\ \text{avec}\quad\pi=dx/2y\qquad,\tag*{(2.1)}
\]
\[
c_4=-2^4\cdot3a_4,\quad c_6=-2^5\cdot3^3\,a_6,\quad\Delta=2^4(-4a_4^3-27a_6^2)\quad .\tag*{(2.2)}
\]
Posant $Y=2y$, $X=x$, on retombe sur la forme de Weierstrass
\[
Y^2=4X^3-g_2X-g_3\quad,\ \text{avec}\quad\pi=d\,X/Y\quad,\ \text{et}\tag*{(2.3)}
\]
\[
g_2=-4a_4=\frac{1}{12}c_4\qquad g_3=-4a_6=\frac{1}{216}c_6\tag*{(2.4)}
\]
\[
\Delta=g_2^3-27g_3^2\qquad .
\]

\uline{Proposition} 2.5.\qquad \uline{Au dessus de} $\mathrm{Spec}(\mathbb Z[2^{-1},3^{-1}]$, \uline{il existe un schéma de modules pour les courbes de genre un munies d'une forme différentielle invariante inversible}. \uline{C\unsure{;|e} schéma est}
\[
\mathrm{Spec}(\mathbb Z[2^{-1},3^{-1}][g_2,g_3]\qquad ,
\]
% archive sheet number: 8
\scanpage{10}
\mspage{7}
\uline{avec pour courbe universelle} (\uline{en coordonnées non homogènes}) \handwritten{\tikz[baseline=-.2ex,x=1ex,y=1ex]{\draw (0,0)--(.32,1.18) .. controls (.15,1.27) and (.46,1.55) .. (.54,1.38) .. controls (.62,1.27) and (.40,1.13) .. (.32,1.18);}}\inserted{\handwritten{section $e$}}
\[
Y^2=4X^3-g_2X-g_3
\]
\uline{et pour différentielle invariante} $\pi=d\,X/Y$.

De plus, lorsque 2 et 3 sont inversibles, la donnée d'une forme différentielle invariante suffit à rigidifier une courbe de genre 1.

Il résulte aussitôt de (2.2) qu'on a
\[
c_4^3-c_6^2=1728\,\Delta\tag*{(2.6)}
\]
et cette formule, se réduisant à une identité algébrique, est valable en toute caractéristique.

3.\qquad \uline{Cas où} 2 \uline{et} $c_4$ \uline{sont inversibles}.

Supposons que $E$ soit une courbe de genre un sur une base $S$ où 2 soit inversible. La forme différentielle invariante $\pi$ étant donnée, on peut choisir $x$ et $y$ tels que dans (1.3) on ait $a_1=a_3=0$. Les changements de coordonnées respectant cette condition s'obtiennent en faisant $u=1$, $s=t=0$ dans (1.2).

On a donc
\[
y^2=x^3+a_2x^2+a_4x+a_6\tag*{(3.1)}
\]
\[
b_2=4a_2\qquad b_4=2a_4\qquad b_6=4a_6\qquad d_6=8(a_2a_4-9a_6)
\]
\[
c_4=2^4(a_2^2-3a_4)\qquad -c_6=2^6a_2^3-2^5\,3^2\,a_2a_4+2^5\,3^3\,a_6\quad\text{et}\tag*{(3.2)}
\]
\[
\Delta=2^4\,\delta \qquad ,
\]
% archive sheet number: 9
\scanpage{11}
\mspage{8}
où $\delta$ est le discriminant du second membre de (3.1). Le discriminant $\delta$ étant invariant par translation, il suffit de vérifier cette identité algébrique pour 3 inversible et $a_2=0$.

Lorsque $c_4$ est inversible, un et un seul choix de $x$ et $y$ donne $d_6=0$ \unsure{?|,} i.e. $a_2a_4-9a_6=0$. On a, lorsque cette condition est remplie,
\[
-c_6=2^6(a_2^3-3^3a_6)\qquad \Delta=2^6(-a_2^3a_6-a_4^3+2\cdot3^3\,a_6^2)\tag*{(3.3)}
\]
et
\[
\begin{aligned}
a_2c_4&=-2^{-2}c_6\\
a_4c_4^2&=-2^2\,3^2\,\Delta\\
a_6c_4^3&=c_6\,\Delta\qquad .
\end{aligned}\tag*{(3.4)}
\]

Il suffit de vérifier ces identités en caractéristique 0 , auquel cas la troisième est le produit des deux premières, puisque $d_6=0$.

\uline{Proposition} 3.5.\qquad \uline{Au dessus de} $\mathrm{Spec}(\mathbb Z[2^{-1}])$, \uline{il existe un schéma de module pour les courbes de genre un munies d'une forme différentielle invariante inversible et telles que} $c_4$ \uline{soit inversible. Ce schéma est}
\[
\mathrm{Spec}(\mathbb Z[2^{-1}][a_2,a_4,a_6,(a_2^2-3a_4)^{-1}]/(a_2a_4-9a_6))
\]
\uline{avec pour courbe universelle} (\uline{en coordonnées non homogène}) \handwritten{\tikz[baseline=-.2ex,x=1ex,y=1ex]{\draw (0,0)--(.26,1.20) .. controls (.13,1.47) and (.55,1.58) .. (.47,1.34) .. controls (.42,1.20) and (.29,1.16) .. (.26,1.20);}}\inserted{\handwritten{section $e$}}
\[
y^2=x^3+a_2x^2+a_4x+a_6
\]
% archive sheet number: 10
\scanpage{12}
\mspage{9}
\uline{et pour différentielle invariante} $\pi=dx/2y$.

De plus, sous ces hypothèses, la donnée d'une forme différentielle invariante suffit à rigidifier une courbe.

4.\qquad \uline{Cas où} 2 \uline{est nilpotent et} $c_4$ \uline{inversible}.

Lorsque $2=0$, la forme différentielle ne suffit plus jamais à rigidifier une courbe de genre un. Lorsque $c_4$ est inversible et 2 nilpotent, donc $a_1$ inversible, en changeant le choix de $\pi$, $x$, $y$ on peut trouver une équation
\[
y^2+xy=x^3+a_2x^2+a_6\qquad \pi=\frac{dx}{x+2y}\qquad .\tag*{(4.1)}
\]
Les autres choix de $x$ et $y$ donnant lieu à une telle équation s'obtiennent en faisant, dans (1.2)
\[
r=t=0\qquad u=1+2s\tag*{(4.2)}
\]
et on aura
\[
\begin{aligned}
u^2a'_2&=a_2-s-s^2\\
u^6a'_6&=a_6
\end{aligned}\tag*{(4.3)}
\]
et \uline{modulo 2}
\[
c_4=1\qquad c_6=1\qquad \Delta=a_6\qquad .\tag*{(4.4)}
\]

5.\qquad \uline{Cas où} $c_4=0$, $p=2$ \uline{ou} $p=3$

a) $p=3$

En carac\deleted{\unsure{e|b|t}}téristique 3 , lorsque l'équation est mise sour forme

% archive sheet number: 11
\scanpage{13}
\mspage{10}
(3.1), $c_4=0$ si et seulement si $a_2^2=0$. Supposons qu'on ait même $a_2=0$
\[
y^2=x^3+a_4x+a_6\qquad \pi=a_4^{-1}\,dy=dx/2y\tag*{(5.1)}
\]
\[
c_4=c_6=0\qquad \Delta=-a_4^3\qquad .
\]

Les autres systèmes $(\pi,x,y)$ donnant une équation de cette forme s'obtiennent en faisant $s=t=0$ dans (1.2) et
\[
\begin{aligned}
u^4a'_4&=a_4\\
u^6a'_6&=a_6+ra_4+r^3\qquad .
\end{aligned}\tag*{(5.2)}
\]

b) $p=2$

En caractéristique 2 , on a $c_4=a_1^4$. Supposons qu'on ait $a_1=0$ ; $\pi$ étant donnée, on peut alors choisir $x$ et $y$ tels que (1.3) soit
\[
y^2+a_3y=x^3+a_4x+a_6\qquad\pi=\frac{dx}{a_3}=dy/\raisebox{-.45ex}{$x^2+a_4$}\unsure{:|!}\tag*{(5.3)}
\]
\[
b_2=b_4=0\qquad b_6=a_3^2\qquad b_8=a_4^2\qquad\Delta=a_3^4\qquad c_4=c_6=0\, .
\]
Les autres systèmes $(\pi,x,y)$ donnant une équation de cette forme s'obtiennent en faisant $r=s^2$ dans (1.2) et
\[
\begin{aligned}
u^3\unsure{a}'_3&=a_3\\
u^4a'_4&=a_4+sa_3+s^4\\
u^6a'_6&=a_6\unsure{-|\cdot}+s^2a_4+ta_3+s^6+t^2\qquad .
\end{aligned}\tag*{(5.4)}
\]
% archive sheet number: 12
\scanpage{14}
\mspage{11}
6.\qquad \uline{Applications}.

\uline{Proposition} 6.1.\qquad \uline{Si} $E$ \uline{est une courbe de genre un sur un corps algébriquement clos}, $E$ \uline{est de type elliptique} (\uline{resp. multiplicatif}, \uline{resp. additif}) \uline{si et seulement si} $\Delta\ne0$ (\uline{resp.} $\Delta=0$ et $c_4\ne0$, \uline{resp.} $\Delta=c_4=0$).

a)\qquad Lorsque la caractéristique $p\ne2$, la première assertion résulte de la troisième formule (3.2). Lorsque $p=2$ et $c_4\ne0$, on vérifie aussitôt que la différentielle de (4.1) ne s'annule qu'en $(0,0)$, et $\Delta$ est elliptique si et seulement si $\Delta=a_6\ne0$. Pour $p=2$ et $c_4=0$, on raisonne de même sur (6.2).

b)\qquad Lorsque $p\ne2$, il résulte de (3.2) que $\Delta=c_4=0$ si et seulement si le deuxième membre de (3.1) a toutes ses racines égales, ce qui caractérise le cas additif.

c)\qquad Lorsque $p=2$, $\Delta=0$, $c_4\ne0$, la partie principale de (4.1) en $(0,0)$ est $y^2+xy+a_2x^2=0$, ce qui correspond à un point double à tangentes distinctes.

d)\qquad Lorsque $p=2$, $\Delta=0$, $c_{4\unsure{.}}=0$, la partie principale de (6.2) en $(0,0)$ est $y^2+a_4x+a_6=0$, ce qui correspond à une singularité cuspidale.

Soit alors $E$ une courbe de genre un sur une base $S$, sans fibre de type additif. Puisque $\Delta$ et $c_4$ ne sont jamais simultanément nuls, et sont respectivement des sections de $\omega^{\otimes12}$ et $\omega^{\otimes4}$, on définit une section $j$ de $P_1/S$ par la formule
% archive sheet number: 13
\scanpage{15}
\mspage{12}
\[
j=c_4^3/\Delta=1728+c_6^2/\Delta\quad .\tag*{(6.2)}
\]

\uline{Proposition} 6.3.\qquad \uline{Soient} $E$ \uline{et} $E'$ \uline{deux courbes de genre un sur} $S$, \uline{sans fibre additive}.

(I) $\underline{\mathrm{Isom}}_{\handwritten{S}}(E,E')$ \uline{est représentable, fini et non ramifié sur} $S$

(II) \uline{si} $j_E=j_F$, \uline{la projection de} $\underline{\mathrm{Isom}}_{\handwritten{S}}(\unsure{E|\mathbb E},E')$ \uline{sur} $S$ \uline{est surjective}

(III) \uline{si} $j_E=j_F$ \uline{ne prend pas les valeurs} 0 \uline{et} 1728, $\underline{\mathrm{Isom}}_{\handwritten{S}}(E,E')$ \uline{est un revêtement étale de rang} 2 \uline{de} $S$.

Donnons-nous $E$ et $E'$ par des équations (1.3) et (1.3)'. Se donner un isomorphisme entre $E$ et $E'$ revient à se donner $u$, $r$, $s$ et $t$, vérifiant les équations (1.6), donc aussi (1.7) et (1.8), $u$ devant de plus être inversible. Par hypothèse, localement sur $S$, $c_4$ ou $\Delta$ est inversible (6.1), et de même $c'_4$ ou $\Delta'$ l'est. Localement sur $S$, les équations (1.8) sont donc des équations de dépendance intégrale pour $u$ et $u^{-1}$, les équations $(1.7)_6$ ou $(1.7)_8$ des équations de dépendance intégrale pour $r$ et enfin $(1.6)_6$ (resp. $(1.6)_2$) est une équation de dépendance intégrale pour $t$ (resp. $s$). Ceci prouve que $\underline{\mathrm{Isom}}(E,E')$ est représentable par un schéma affine fini sur $S$ ; il est même non ramifié car les courbes considérées n'ont pas d'automorphismes infinitésimaux (respectant l'origine).

Lorsque 2 est inversible sur $S$, l'assertion (III) résulte de l'énoncé plus précis.

\uline{Proposition} 6.4.\qquad \uline{Soient} $E$ \uline{et} $E'$ \uline{deux courbes de genre un sur} $S$,
% archive sheet number: 14
\scanpage{16}
\mspage{13}
\uline{sans fibre additive}, \uline{munies de formes différentielles invariantes} $\pi$ \uline{et} $\pi'$, \uline{et de même invariant modulaire}. \uline{On suppose que} 2 , $c_4$ \uline{et} $c_6$ \uline{sont inversibles}, \uline{i.e.} \uline{que} $2^{-1}$, $1/j$ et $1/(j-1728)$ \uline{sont des sections de} $S$. \uline{Alors}, \uline{le schéma} $\underline{\mathrm{Isom}}(E,F)$ \uline{s'identifie au schéma défini par l'équation}
\[
\lambda^2=c_4c'_6/c_6c'_4\qquad .
\]

En vertu de (3.6), se donner $(E,\pi)$ revient à se donner $a_2$, $a_4$ et $a_6$ vérifiant $a_2a_4-9a_6=0$, et un isomorphisme de $E$ sur $E'$ est décrit par $\lambda$ tel que $(E\,\pi)$ soit isomorphe à $(E',\lambda\,\pi')$, i.e. par $\lambda$ vérifiant
\[
a'_2=\lambda^2a_2\qquad a'_4=\lambda^4a_4\qquad a'_6=\lambda^6a_6\qquad .\tag*{(6.4.1)}
\]

Par hypothèse, $j=j'$, de sorte que $(c_4^3,c_6^2,\Delta)$ est proportionnel à $({c'_4}^3,{c'_6}^2,\Delta')$, et que $\mu=c_4c'_6/\raisebox{-.45ex}{$c_6c'_4$}$ vérifie
\[
c'_4=\mu^2c_4\qquad c'_6=\mu^3c_6\qquad \Delta'=\mu^6\,\Delta\qquad .\tag*{(6.4.2)}
\]
En vertu de (3.4), on a alors\margin{\handwritten{\unsure{--}}}
\[
a'_2=\mu\,a_2\qquad a'_4=\mu^2a_4\qquad a'_6=\mu^3a_6\qquad .
\]
De plus, d'après (3.4) encore, localement, soit $a_2$, soit $a_4$ et $a_6$ sont inversibles, de sorte que (6.4.1) équivaut à $\lambda^2=\mu$.

Pour prouver (III) en général, on peut se ramener à supposer $S$ noethérien, et supposer soit 2 inversible (cas traité en (6.4)), soit
% archive sheet number: 15
